\section{Construindo, tagueando e publicando imagens}

\subsection{Building images}

O comando mais básico para construir uma imagem é:

\begin{verbatim}
docker build .
\end{verbatim}

O \texttt{.} indica o \textit{build context} --- o diretório onde o builder encontrará o Dockerfile e os arquivos referenciados. Ao executar, o builder baixa a imagem base (se necessário) e executa as instruções do Dockerfile. Sem uma tag, a imagem recebe apenas um ID (hash SHA256), que não é fácil de lembrar.

\subsection{Tagging images}

Tags dão nomes memoráveis às imagens. A estrutura completa de um nome de imagem é:

\begin{verbatim}
[HOST[:PORT_NUMBER]/]PATH[:TAG]
\end{verbatim}

\begin{itemize}
    \item \textbf{HOST}: registro onde a imagem está. Padrão: \texttt{docker.io}.
    \item \textbf{PORT\_NUMBER}: porta do registro (se um host foi especificado).
    \item \textbf{PATH}: caminho no formato \texttt{[NAMESPACE/]REPOSITORY}. Se nenhum namespace é dado, usa-se \texttt{library}.
    \item \textbf{TAG}: identificador de versão. Padrão: \texttt{latest}.
\end{itemize}

Para dar tag durante o build, usa-se a flag \texttt{-t}:

\begin{verbatim}
docker build -t meu-usuario/minha-imagem .
\end{verbatim}

Para adicionar outra tag a uma imagem já existente:

\begin{verbatim}
docker image tag meu-usuario/minha-imagem outro-usuario/outra-imagem:v1
\end{verbatim}

\subsection{Publishing images}

Para enviar a imagem ao registro, é necessário estar autenticado e usar o comando \texttt{docker push}:

\begin{verbatim}
docker login
docker push meu-usuario/minha-imagem
\end{verbatim}

\subsection{Docker Hub}

Para compartilhar Docker images, é necessário um local para armazená-las. O Docker Hub é o registro padrão, provendo tanto um local para armazenar images próprias quanto para encontrar images de outros usuários.

Ao escolher images, o Docker Hub fornece categorias de images confiáveis:

\begin{itemize}
    \item \textbf{Docker Official Images (DOI)}: images selecionadas dos softwares mais populares, seguindo boas práticas e atualizadas regularmente.
    \item \textbf{Docker Hardened Images (DHI)}: images desenhadas para reduzir a superfície de ataque --- seguras, simples e praticamente sem vulnerabilidades.
\end{itemize}

\subsection{Guia prático completo}

\begin{verbatim}
# Clonar o repositório
git clone https://github.com/docker/getting-started-todo-app
cd getting-started-todo-app

# Construir a imagem
docker build -t SEU_USERNAME/concepts-build-image-demo .

# Verificar a imagem criada
docker image ls

# Visualizar histórico de layers
docker image history SEU_USERNAME/concepts-build-image-demo

# Enviar ao Docker Hub
docker push SEU_USERNAME/concepts-build-image-demo
\end{verbatim}

O comando \texttt{docker image history} mostra as camadas da imagem, distinguindo as que foram adicionadas pelo usuário das herdadas da imagem base. Se ao fazer push receber o erro \texttt{requested access to the resource is denied}, verifique se está logado (\texttt{docker login}) e se o username na tag está correto.
